\subsection{Uniform cancellation between insertion and the corrected marginal}\label{sec:B4T09}
Retain the hypotheses, normalizations, and small branch of Sections~\ref{sec:B4T01}--\ref{sec:B4T08}. The parent correction is $\delta_c=(R\nabla\ell)_c/p_c$.
We do not identify the corrected marginal with the true marginal of the count.

\paragraph{Derivative of the prefactor in the variables of the unit-diagonal matrix.}
The exact identity of Section~\ref{sec:B4T03} and the resolvent identity for $S=S_0+E$ give
\[
 G_{ab}-\Gamma_{ab}=-3m_{2,a}\mathbf1_{a=b}+O(k^{-2}).
\]
Indeed, $E_{aa}=3m_{2,a}+O(k^{-2})$, $E_{ab}=O(k^{-2})$ off the diagonal, $\|E\|_\infty=O(k^{-1})$, $\Gamma=I+O_{\max}(k^{-1})$, and $\|\Gamma\|_\infty=O(1)$.
Expanding the exact formula for $g=\nabla\ell$ gives
\[
 g_i+1=-6p_i-f_i+\xi_{r_i}+\zeta_{c_i}+O(k^{-2}),
\]
where we may take
\[
 \xi_r=-\tfrac12(\Gamma_{rr}-1)+\tfrac32m_{2,r},\qquad
 \zeta_c=-\tfrac12(\Gamma_{cc}-1)+\tfrac32m_{2,c}.
\]
Applying the tangent response annihilates the RC functions.
Also, with $R_0=P_p-P_pA\Gamma A^TP_p$, we have $\|P_p^{-1}(R-R_0)\|_\infty=O(k^{-1})$ by the resolvent identity.
Since $-6p-f=O(k^{-1})$, we may replace $R$ by $R_0$ with relative accuracy $O(k^{-2})$. Therefore
\[
 \delta_c=-6p_c-f_c+
 \sum_i p_i(a_i^T\Gamma a_c)(6p_i+f_i)+O(k^{-2}).
\]
On $F\setminus c$, we have $a_i^T\Gamma a_c=1+O(k^{-1})$.
On $E'$, Section~\ref{sec:B4T06} gives $p_i a_i^T\Gamma a_c=-y_i+O(p_i/k^2)$.
The exceptional term at $c$ changes the result by only $O(k^{-2})$, so
\[
 \boxed{\delta_c=-6p_c-f_c+
 \sum_Fp_i(6p_i+f_i)-\sum_{E'}y_i(6p_i+f_i)+O(k^{-2}).}
\]
All sums of errors take into account the mass bounds $\sum_Ep=k$ and $\sum_Fp=O(1)$.

\paragraph{Comparison with the child.}
Let $\epsilon_c=L(B+c)-L(B)-\log p_c$. Combining the energy from Section~\ref{sec:B4T06}, the substitution from Section~\ref{sec:B4T07}, and the determinant expansion from Section~\ref{sec:B4T08} gives
\[
 \begin{aligned}
 \epsilon_c={}&-6p_c-f_c+6\sum_Fp_i^2-6\sum_{E'}p_iy_i
 -\sum_{E'}y_i f_i\\
 &+\sum_Fp_i f_i^I-\sum_{j\in J}(K_0)_{jj}-1/k+O(k^{-2}).
 \end{aligned}
\]
The $\tfrac12\log(2\pi)$ terms agree because both supports have connected RC graphs. The cofactors already account for the covolume.
Subtracting the formula for $\delta_c$ leaves
\[
 \epsilon_c-\delta_c=-\sum_{j\in J}(K_0)_{jj}-1/k
                    -\sum_Fp_i(f_i-f_i^I)+O(k^{-2})=O(k^{-2})
\]
by the final identity of Section~\ref{sec:B4T08}.
Since $\delta_c=O(k^{-1})$,
\[
 \boxed{r_c=L(B+c)-L(B)-\log[p_c(1+\delta_c)]=O(k^{-2}).}
\]
The constants depend only on $\delta_0,C,\eta,\gamma$, and there exists $k_0(\delta_0,C,\eta,\gamma)$ such that the result holds for every $k\ge k_0$ and every legal child of the parent, without selecting children.

\paragraph{Scope.}
The child stationary point used here lies on the branch $p<1/4$. Section~\ref{sec:B4T02} proves existence and uniqueness on this branch.
We do not assert that the quantitative hypotheses hold for every residual board in the counting problem.
